
\documentclass{article}
\usepackage{amsmath}
\usepackage{graphicx}
\usepackage{enumitem}

\title{Empirical Exercises for Economics of Global Shifts Course}
\author{}
\date{}

\begin{document}

\maketitle

\section*{1. Global Trade and Economic Growth}
\textbf{Objective:} Analyze the relationship between global trade and economic growth in different countries.

\textbf{Instructions:}
\begin{itemize}
    \item Collect data on GDP and trade volume (exports + imports) for a selection of countries over the past 30 years (you can use databases like the World Bank, IMF, or UNCTAD).
    \item Plot the trade-to-GDP ratio and growth rates for each country over time.
    \item Perform a correlation analysis between trade volume and GDP growth for each country.
    \item Discuss the results: Do higher levels of trade seem to correlate with higher economic growth? Are there differences between developed and developing countries?
\end{itemize}

\textbf{Empirical Tasks:}
\begin{enumerate}
    \item Compute the trade-to-GDP ratio for each country.
    \item Run a regression analysis to determine the relationship between trade and GDP growth.
    \item Evaluate whether there are time-lags between trade expansion and GDP growth.
\end{enumerate}

\section*{2. Technological Advancements and Productivity Growth}
\textbf{Objective:} Investigate the role of technological advancements (e.g., ICT, AI) in productivity growth.

\textbf{Instructions:}
\begin{itemize}
    \item Use a dataset such as the OECD’s productivity database or World Bank data on labor productivity and technological innovation indicators.
    \item Choose a few countries (e.g., the U.S., China, India, and Germany) and gather data on productivity (output per hour worked) and technological investments (R\&D expenditure, patents granted).
    \item Analyze whether higher levels of technological innovation are associated with higher productivity growth.
\end{itemize}

\textbf{Empirical Tasks:}
\begin{enumerate}
    \item Calculate labor productivity growth rates for each country over the last 10 years.
    \item Run a regression analysis to see if there is a statistically significant relationship between technological innovation (measured by R\&D spending or patents) and productivity growth.
    \item Compare the results across countries with different levels of technological development.
\end{enumerate}

\section*{3. Automation and Employment}
\textbf{Objective:} Examine the impact of automation on labor markets, especially in manufacturing sectors.

\textbf{Instructions:}
\begin{itemize}
    \item Collect data on automation adoption rates in key industries (e.g., manufacturing, retail) and employment data (total jobs and sector-specific jobs) from the Bureau of Labor Statistics (U.S.) or national labor data sources.
    \item Focus on countries with varying levels of automation, such as the U.S., South Korea, and Germany.
    \item Analyze the impact of automation on employment trends: Are jobs being displaced, or is there a shift to higher-skill employment?
\end{itemize}

\textbf{Empirical Tasks:}
\begin{enumerate}
    \item Calculate the change in employment levels in automation-heavy industries over the last 20 years.
    \item Use a difference-in-differences (DiD) approach to compare regions with varying automation adoption rates.
    \item Interpret the impact of automation on job displacement and wage levels in manufacturing versus service sectors.
\end{enumerate}

\section*{4. Foreign Direct Investment (FDI) and Economic Growth}
\textbf{Objective:} Investigate the role of foreign direct investment (FDI) in stimulating economic growth in developing countries.

\textbf{Instructions:}
\begin{itemize}
    \item Use the World Bank or UNCTAD database to gather data on FDI inflows and GDP growth for a set of developing countries (e.g., Brazil, India, South Africa).
    \item Examine how changes in FDI inflows are correlated with GDP growth and other economic indicators such as employment and infrastructure development.
\end{itemize}

\textbf{Empirical Tasks:}
\begin{enumerate}
    \item Create a time series plot of FDI inflows and GDP growth for each country.
    \item Perform a regression analysis to see if there is a statistically significant impact of FDI on GDP growth.
    \item Compare the growth rates of countries with high FDI inflows to those with low FDI inflows and discuss the potential reasons for observed differences.
\end{enumerate}

\section*{5. Digital Economy and Inequality}
\textbf{Objective:} Assess the effect of the digital economy on income inequality across countries.

\textbf{Instructions:}
\begin{itemize}
    \item Collect data on internet penetration rates, digital services usage (e.g., e-commerce, digital banking), and income inequality (Gini coefficient) for a set of countries over the last 20 years.
    \item Investigate whether there is a relationship between the expansion of digital services and changes in income inequality.
\end{itemize}

\textbf{Empirical Tasks:}
\begin{enumerate}
    \item Compute the Gini coefficient for a selection of countries over time and compare it with the level of digital penetration.
    \item Run a regression model to test whether increased digital services usage is associated with lower or higher income inequality.
    \item Discuss the potential channels through which the digital economy might affect inequality (e.g., access to education, job creation, wealth accumulation).
\end{enumerate}

\section*{6. Global Supply Chains and Manufacturing Relocation}
\textbf{Objective:} Explore the effects of globalization and global supply chains on manufacturing locations and employment.

\textbf{Instructions:}
\begin{itemize}
    \item Use data on global manufacturing output, trade flows, and employment (e.g., from the WTO or World Bank).
    \item Examine the movement of manufacturing jobs from high-wage countries to low-wage countries over the past few decades (e.g., U.S. to China, Mexico to Central America).
    \item Analyze the effects on wages, employment rates, and regional economic development.
\end{itemize}

\textbf{Empirical Tasks:}
\begin{enumerate}
    \item Plot trends in manufacturing output and employment in major countries over time.
    \item Use a panel data regression model to assess how changes in global supply chain integration (e.g., outsourcing, offshoring) correlate with domestic manufacturing employment and wages.
    \item Evaluate the impact on income inequality within countries affected by offshoring.
\end{enumerate}

\section*{7. Environmental Impact of Technological Change}
\textbf{Objective:} Analyze the environmental consequences of technological shifts in key industries.

\textbf{Instructions:}
\begin{itemize}
    \item Collect data on energy consumption, carbon emissions, and technological advancements in energy efficiency or green technologies.
    \item Focus on industries such as energy production, manufacturing, and transportation.
    \item Examine the environmental effects of technological advancements in these sectors and their relationship to economic growth.
\end{itemize}

\textbf{Empirical Tasks:}
\begin{enumerate}
    \item Analyze the relationship between technological innovation in renewable energy (e.g., wind, solar) and reductions in carbon emissions.
    \item Calculate the carbon intensity of GDP (CO2 emissions per unit of GDP) for countries with varying levels of technological adoption.
    \item Perform a regression analysis to determine the effect of green technologies on reducing environmental degradation without sacrificing economic growth.
\end{enumerate}

\section*{8. Demographic Shifts and Economic Outcomes}
\textbf{Objective:} Investigate how demographic changes (e.g., aging populations) influence economic outcomes such as growth, savings rates, and labor force participation.

\textbf{Instructions:}
\begin{itemize}
    \item Use demographic data from the UN Population Division and economic data (GDP growth, savings rates, labor force participation) from the World Bank.
    \item Focus on countries with rapidly aging populations (e.g., Japan, Germany) and compare them to countries with younger populations (e.g., India, Brazil).
\end{itemize}

\textbf{Empirical Tasks:}
\begin{enumerate}
    \item Analyze the relationship between the aging population and GDP growth rates across countries.
    \item Perform a regression analysis to assess the impact of demographic aging on national savings rates and labor force participation.
    \item Discuss potential policy responses to mitigate the economic effects of demographic shifts, such as pension reform or changes in immigration policy.
\end{enumerate}

\end{document}
